\chapter{From economic feedback to a research question}\label{ch:research}

\section{A gauge is useful only in an arrangement}
Vale installs a clear gauge on the clinic's tank. The gauge makes the level visible. For several days, everyone admires the improvement. Then the level falls and nobody opens the valve, because nobody has been assigned responsibility. The gauge has worked perfectly. The arrangement has not provided water.

Now imagine the opposite case. A caretaker opens the valve whenever a neighbour asks, but cannot see the tank level and has no reliable record of what arrived. The arrangement contains action and concern, yet lacks essential feedback. A third case is also possible: the level is visible and responsibility is assigned, but the pump has no power. Information, decision and physical capability are distinct requirements.

These simple cases help clarify the claim that an economy needs a connection to physical functioning. The claim should not be expressed as though economies contain no feedback. Prices, inventories, contracts, public decisions and engineering controls can all participate in feedback arrangements. The research question is whether selected physical consequences are observed, evaluated and acted upon adequately for a specified purpose.

EBU initially addresses one part of this arrangement: an explicit signed account of a represented physical action. It does not become a controller merely because its result can inform a controller. The distinction will prevent us from attributing behaviour to an equation that only evaluates a change.

\section{Five questions behind a feedback loop}
The first question is what is being observed. A gauge may show stock, pressure or flow. Those are different quantities. A meter showing yesterday's delivery cannot directly tell the caretaker how much is in the tank now unless withdrawals and other changes are also known. Observations need a declared meaning and time.

The second question is what reference is used. A caretaker cannot decide whether a reading calls for action without some criterion. The criterion might concern an intended stock, a service schedule or a safety condition. In our later Gaussian account, a declared reference and positive scale determine the represented evaluation. They do not establish the complete service policy.

The third question is what action is available. A valve can transfer water only if an admissible route and sufficient source stock exist. A displayed shortage does not create a pipe. The set of physically possible actions therefore constrains every decision rule, including any rule that uses EBU.

The fourth question is how an action is selected. A person, protocol or controller may choose among available actions. Selection might prioritise service, reduce a declared potential or follow another justified rule. Different selectors can use the same physical account and produce different behaviour. Their differences must be tested rather than erased by giving them a common label.

The fifth question is how the result is verified. The requested action may differ from the accepted action, and the accepted action may differ from what was actually performed. A reliable record distinguishes these stages. If a request for four units produces a measured transfer of three, the final physical account should not retain four simply because that was the plan.

\begin{intuitionbox}{The account and the action rule}
A potential evaluates a represented state. An action account compares represented states. A selector chooses an action from an allowed set. Observations establish what happened. These components can work together, but specifying one does not silently specify the others.
\end{intuitionbox}

\section{Turning an aspiration into a question}
``Make the economy support life'' expresses a purpose, but is too broad to be a single scientific hypothesis. To investigate it, we must identify a smaller claim. For example, one could ask whether a declared action rule maintains specified services in a particular resource world under declared disturbances. Such a question needs explicit variables, criteria, alternatives and conditions of failure.

The work in this book comes before that complete study. It asks whether the proposed physical account can be defined consistently and what identities follow from the definition. Can a finite action be evaluated without replacing it with a tangent? Can a local computation recover the affected global difference? Can a sequence of records close when external changes are recorded separately? These are bounded mathematical questions.

The fact that an identity can be proved does not establish that a decision rule using it will achieve the broader purpose. A speedometer can measure speed accurately while a driver takes the wrong road. Likewise, an exact EBU quote can be attached to an action selected by a poor policy. Accuracy of the account and adequacy of the action rule must each be examined.

Conversely, a successful-looking trajectory does not repair an inconsistent account. If a method double-counts resource creation or attributes an external change to an actor, a desirable endpoint does not make the accounting correct. The programme needs both a trustworthy representation and a fair test of behaviour.

\section{The four kinds of statement}
A \emph{definition} declares what a symbol or model object means. Defining a potential does not demonstrate that it measures every important burden. It makes the subsequent mathematical question precise. The definition can be useful, inadequate or inappropriate for a proposed application.

A \emph{theorem} establishes a conclusion from specified assumptions. If a fixed potential is evaluated before and after a sequence of actions, certain sums may telescope exactly. The theorem remains conditional on the stated sequence, representation and accounting rules. Removing an assumption requires a new argument.

An \emph{implementation check} establishes something about a program at the checked boundary. A test that verifies an arithmetic identity on a fixture is evidence about the implementation of that identity for those conditions. It is not an experiment establishing long-run homeostasis, and the number of passing checks does not change the category.

An \emph{empirical finding} records what an executed, adequately documented investigation observed. Its interpretation depends on the design, measurements, comparisons and declared limitations. A result in one world is not automatically a result in another. The distinction matters even when the same code and the same mathematical notation appear in both.

A candidate design belongs to none of these categories by wish alone. It may contain definitions and proposed procedures, but it remains prospective until the required adoption, implementation and execution occur. This vocabulary will appear throughout the book because the status of a claim is part of its meaning.

\section{The temptation to build a favourable world}
Suppose we wish to compare two hypothetical rules. One transfers toward a declared stock reference. Another serves requests in their arrival order. It would be easy to choose examples in which every request happens to move stock toward the reference, then announce agreement between EBU guidance and service. It would be equally easy to choose only examples in which the request rule performs badly.

A useful investigation would include cases that separate the rules. A necessary service may move stock away from a reference. A reference-restoring action may provide no requested service. A source-preserving rule may refuse all transfers. These cases expose what the rules actually prioritise.

The reference itself must also be declared before inspecting the outcomes that are meant to test it. Otherwise, an evaluator can move the target after seeing where the favoured rule went. An action that previously increased burden can be made to appear helpful by redefining the reference around its endpoint. That is a change in evaluation, not a discovery about the original action.

Examples in a textbook serve a different purpose from registered trials: they illustrate a distinction openly. The Ridge--Vale quantities are chosen to make arithmetic legible. We may change them to teach a different point, provided the change is stated. We may not present the resulting teaching convenience as evidence that an untested controller succeeds generally.

\section{What the first physical model includes}
The core mathematical world uses a finite collection of places and one homogeneous resource. Stocks are nonnegative, and internal lossless transfers preserve a fixed total. The declared reference is compatible with that total. Positive fixed scales determine how deviations from the reference are measured. These assumptions are strong enough to support exact calculations and narrow enough to remain visible.

The world does not acquire regeneration because Ridge is called a source in a story. A source of an individual transfer is simply the place from which that transfer withdraws stock. If natural production, consumption or disturbance is introduced, its physical effect needs its own declaration and record. It must not be treated as an actor's contribution merely because it improves the state.

Nor does the Gaussian potential create movement. A state away from the reference can remain there if no action occurs. A graph can make some desired redistributions unavailable. A permission rule can limit an otherwise attractive transfer. These possibilities are not edge cases added to spoil an elegant theory; they are consequences of distinguishing state evaluation from dynamics.

The model's simplicity also helps us see a basic problem in some proposed demonstrations. If a world starts exactly at its reference and contains no disturbance or drive, a rule that acts only on a restoring field may have no reason to move. That does not test recovery from disturbance or sustained service. The scientific question must fit the declared world.

\section{A worked comparison with no winner declared}
Consider two proposals from the same stock state, $(18,2)$. Proposal A transfers four units, producing $(14,6)$. Proposal B transfers eight units, producing $(10,10)$. Both preserve the total twenty. Suppose the reference is $(10,10)$ and the scales will later be declared equally. Proposal B reaches that reference; proposal A moves toward it.

This information still does not establish a universal winner. A request may have been for four units. The route may allow only four. A separate requirement may prevent sending eight. The physical action set and service purpose must be included before choosing. Even when both proposals are permitted, a controller's objective must be explicit.

Now add Proposal C: do nothing. It preserves the current total and avoids the immediate work of transfer. If success is defined only as ``no stock becomes negative'', C may appear perfectly successful. If a declared service must be supplied, C may fail. This is why success criteria must include the intended task rather than only the easiest invariant.

The worked comparison teaches an evaluation habit. Record physical validity, represented field change and service outcome separately. Then apply the declared interpretation rule. Do not hide all three inside a single word such as ``good'', and do not change the interpretation after discovering which proposal your preferred method selected.

\section{What would count as progress}
Progress can occur at several levels. Clarifying a unit prevents an invalid calculation. Proving an identity establishes a conditional relation. Correcting code makes implementation more faithful. Designing a fair comparison makes a future result interpretable. A negative experimental result can reveal that an elaborate method adds no useful benefit over a simpler one.

None of these achievements needs to masquerade as another. A mathematical book can be valuable before a long-run experiment is performed. Its value lies in explaining definitions and consequences accurately, including what remains unproved. An experimental programme can then target the remaining questions rather than pretending that they were settled by notation.

The central long-run question is demanding: can physically permitted actions guided by a declared account support the intended functions under specified changing conditions? Answering it would require a plant with those changes, a controller, comparison rules, observations and failure criteria. This edition provides the conceptual and mathematical foundation for asking the question. It does not report that the answer is yes.

\section{A question needs a possible disappointing answer}
Suppose a proposed comparison is written so that EBU is declared successful whenever it improves the potential, but a comparator is declared successful only when it improves the potential and serves every request. The asymmetric criteria have already favoured one method before anything runs. The remedy is to define the relevant claims and apply their criteria consistently to every eligible arm.

A comparison can also become uninformative when every possible result is given a favourable interpretation. If EBU acts, the evaluator praises responsiveness; if it does nothing, the evaluator praises preservation; if service fails, the evaluator calls the world unrealistic after the fact. A useful design states in advance which observations would count against its central claim and why.

That requirement does not mean that every failure disproves every part of the project. A controller may fail while the accounting identity remains correct. An unsuitable world may reveal a design error rather than a universal impossibility. The response should locate the failed proposition precisely. This lets the programme learn without turning every setback into total rejection or every success into total vindication.

The explanatory book has a related obligation. Its worked examples should include limits, negative signs and competing purposes, so that readers can see what the equations actually say. A sequence of exclusively favourable stories would teach advocacy more effectively than it taught the model.

\section{The role of institutions in a physical programme}
Even a successful resource controller would not settle ownership, consent, privacy or access. Measurements might concern people or facilities whose participation requires agreed rules. An incorrect receipt might need correction. A reference could disadvantage a community. A decision could impose burdens on people outside the measured boundary.

These are not peripheral issues to be solved by adding a positive coefficient. They concern who may decide, who can contest a decision and which rights constrain the system. A mathematically explicit account can help expose choices, but cannot legitimate them merely by calculating consistently.

One useful design principle follows: keep the route from observation to consequence inspectable. A person should be able to distinguish a meter reading, a model assumption, an action-selection rule and an institutional sanction. When all four are compressed into one unexplained score, neither scientific review nor public accountability can identify where an error entered.

\section{Guided problems and answers}
\subsection{The perfect gauge}
The tank gauge is accurate, but the clinic repeatedly runs out of water. Has the gauge failed?

The observations alone do not show a measurement failure. Responsibility, available actions, decision rules or physical supply may have failed. A feedback arrangement requires more than a sensor. The analogous lesson for EBU is that an exact account does not prove effective control.

\subsection{The passing arithmetic check}
A program reproduces an exact potential difference on one hundred synthetic examples. Has it established long-run homeostasis?

No. It has checked a computational property on the declared examples, assuming the test itself is adequate. Long-run behaviour requires a specified evolution and an authorised investigation. Increasing the fixture count does not turn a static arithmetic question into a dynamic service result.

\subsection{The target moved after the run}
A proposed rule performs poorly relative to its original reference. The evaluator chooses a new reference equal to its final state and reports zero final burden. What changed?

The evaluation changed. The original physical history did not. The new account may be mathematically calculable, but it cannot be substituted for the original registered comparison. The change must be disclosed and any new study treated according to its own prospective rules.
